% ============================================================
% 01_resumen.tex — Resumen, Palabras clave, Abstract, Keywords
% ============================================================

\abstractsection{Resumen}

\noindent
El presente trabajo describe el diseño, desarrollo e implementación de NodeAlert IoT,
un sistema embebido de monitoreo ambiental y control autónomo de iluminación orientado
a entornos educativos. El sistema integra un microcontrolador ESP32-C3 programado bajo
ESP-IDF con FreeRTOS, tres sensores conectados por bus \(I^2C\) —dos fotosensores BH1750
y un sensor de temperatura/humedad SHT30—, y un sensor de presencia por radar de
microondas HLK-LD2420 interfaceado por UART. La lógica de control de iluminación
opera de forma completamente autónoma en el firmware, con histéresis de doble umbral
y un antirebote por software de 2\,s, sin depender de conectividad de red. La
telemetría se publica mediante el protocolo MQTT hacia un servidor Mosquitto
desplegado en un contenedor Docker sobre Raspberry Pi. El backend de procesamiento
se implementó en Node-RED con almacenamiento persistente en SQLite y expone una API
REST que alimenta un frontend web desarrollado en React. El sistema satisface los
requisitos mínimos de la cátedra: programación de microcontrolador,
comunicación bidireccional vía MQTT, broker propio
configurado con autenticación, publicadores y suscriptores diferenciados, e interfaz
de usuario con visualización de variables y controles remotos. El prototipo fue
validado en condiciones de laboratorio demostrando operación autónoma estable y
correcta integración de todos los subsistemas.

\vspace{10pt}
\noindent\textbf{Palabras clave:} Internet de las Cosas, ESP32-C3, MQTT, Node-RED,
control de iluminación, sensor de presencia, BH1750, SHT30, LD2420, Docker.

\vspace{20pt}
\hrule
\vspace{10pt}

\abstractsection{Abstract}

\noindent
This paper describes the design, development, and implementation of NodeAlert IoT,
an embedded environmental monitoring and autonomous lighting control system aimed
at educational environments. The system integrates an ESP32-C3 microcontroller
programmed under ESP-IDF with FreeRTOS, three sensors connected via \(I^2C\) bus
—two BH1750 light sensors and an SHT30 temperature/humidity sensor—, and an
HLK-LD2420 microwave radar presence sensor interfaced over UART. The lighting
control logic operates fully autonomously within the firmware, featuring dual-threshold
hysteresis and a 2\,s software debounce, independently of network connectivity.
Telemetry is published over MQTT to a Mosquitto broker running in a Docker container
on a Raspberry Pi. The processing backend is implemented in Node-RED with SQLite
persistent storage and exposes a REST API that feeds a React-based web frontend.
The system fulfills the course's minimum requirements: microcontroller programming,
bidirectional MQTT communication, a self-configured
authenticated broker, differentiated publishers and subscribers, and a user interface
with variable visualization and remote controls. The prototype was validated under
laboratory conditions, demonstrating stable autonomous operation and correct
integration of all subsystems.

\vspace{10pt}
\noindent\textbf{Keywords:} Internet of Things, ESP32-C3, MQTT, Node-RED,
lighting control, presence sensor, BH1750, SHT30, LD2420, Docker.

\vspace{10pt}
\hrule

